% from ega3-1-fr.tex
\oldpage[III]{123}

désignerons par $\widehat{\mathcal F}$ son complété $\mathcal F_{/X'}$ le
long de $X'$ (I, 10.8.4) qui est un $\mathcal O_{\widehat X}$-Module
cohérent (I, 10.8.8).
\end{env}

\begin{env}[4.1.2]
\label{III.4.1.2-fr}
Soit $\mathcal J$ un Idéal cohérent de $\mathcal O_Y$ tel que
$\operatorname{Supp}(\mathcal O_Y/\mathcal J)=Y'$ (I, 5.2.1)~; on sait
(I, 4.4.5) que
\[
\mathcal K=f^*(\mathcal J)\mathcal O_X
\]
est un Idéal cohérent de $\mathcal O_X$ tel que
\[
\operatorname{Supp}(\mathcal O_X/\mathcal K)=X'.
\]
Nous considérerons pour tout $k\geq0$ les $\mathcal O_X$-Modules cohérents
\[
\mathcal F_k
=\mathcal F\otimes_{\mathcal O_Y}(\mathcal O_Y/\mathcal J^{k+1})
=\mathcal F/\mathcal K^{k+1}\mathcal F.
\]
Les $\mathcal O_Y$-Modules $R^nf_*(\mathcal F)$ et
$R^nf_*(\mathcal F_k)$ sont cohérents pour tout $n$
(\hyperref[III.3.2.1-fr]{3.2.1}). Pour tout $k\geq0$ et tout $n$,
l'homomorphisme canonique $\mathcal F\to\mathcal F_k$ définit par
fonctorialité un homomorphisme
\[
R^nf_*(\mathcal F)\longrightarrow R^nf_*(\mathcal F_k).
\tag{4.1.2.1}\label{III.4.1.2.1-fr}
\]
En outre, comme $\mathcal F_k$ est un
$\mathcal O_X/\mathcal K^{k+1}$-Module, $R^nf_*(\mathcal F_k)$ est un
$\mathcal O_Y/\mathcal J^{k+1}$-Module (0, 12.2.1) et on déduit donc de
(\hyperref[III.4.1.2.1-fr]{4.1.2.1}) un homomorphisme
\[
R^nf_*(\mathcal F)\otimes_{\mathcal O_Y}
(\mathcal O_Y/\mathcal J^{k+1})
\longrightarrow R^nf_*(\mathcal F_k).
\tag{4.1.2.2}\label{III.4.1.2.2-fr}
\]
Les deux membres de (\hyperref[III.4.1.2.2-fr]{4.1.2.2}) forment deux
systèmes projectifs, et la limite projective du premier membre n'est autre
que le complété $(R^nf_*(\mathcal F))_{/Y'}$, que nous noterons
$\widehat{R^nf_*(\mathcal F)}$. En outre, il est immédiat que les
homomorphismes (\hyperref[III.4.1.2.2-fr]{4.1.2.2}) forment un système
projectif, d'où par passage à la limite un homomorphisme canonique
\[
\varphi_n:\widehat{R^nf_*(\mathcal F)}
\longrightarrow\varprojlim_kR^nf_*(\mathcal F_k).
\tag{4.1.2.3}\label{III.4.1.2.3-fr}
\]
D'ailleurs (\hyperref[III.4.1.2.2-fr]{4.1.2.2}) est un homomorphisme de
$(\mathcal O_Y/\mathcal J^{k+1})$-Modules, et par suite (I, 10.8.3) peut
être considéré comme un homomorphisme continu de
$\mathcal O_{\widehat Y}$-Modules topologiques pseudo-discrets
(0\textsubscript{I}, 3.8.1). L'homomorphisme $\varphi_n$ est par suite un
homomorphisme continu de $\mathcal O_{\widehat Y}$-Modules topologiques.
\end{env}

\begin{env}[4.1.3]
\label{III.4.1.3-fr}
Soit $i:\widehat X\to X$ le morphisme canonique d'espaces annelés défini
dans (I, 10.8.7), de sorte que l'on a le diagramme commutatif
\[
\xymatrix{
X_k\ar[r]^{h_k}\ar[dr]_{i_k}&\widehat X\ar[d]^i\\
&X
}
\tag{4.1.3.1}\label{III.4.1.3.1-fr}
\]
où $X_k$ est le sous-préschéma fermé de $X$ défini par l'Idéal
$\mathcal K^{k+1}$, $i_k$ l'injection canonique, $h_k$ le morphisme
d'espaces annelés correspondant à l'identité dans les espaces sous-jacents
et à l'homomorphisme canonique
$\mathcal O_{\widehat X}\to\mathcal O_X/\mathcal K^{k+1}$ (I, 10.5.2).
En outre, on a $\widehat{\mathcal F}=i^*(\mathcal F)$ (I, 10.8.8) à un
isomorphisme canonique près. On sait que l'on a
\[
H^n(X_k,i_k^*(\mathcal F_k))=H^n(X,\mathcal F_k)
\tag{4.1.3.2}\label{III.4.1.3.2-fr}
\]

